\section{Extending structures}\label{sec:02-functions-and-structures:03-extending-structures:1}

\subsection{Recall}\label{sec:02-functions-and-structures:03-extending-structures:2}

Formal definitions cited in this section, gathered here for quick reference (full citations in the \hyperref[sec:root:bibliography:1]{Bibliography}):

\begin{itemize}
\tightlist
\item
  \textbf{Forgetful functor.} ``A functor which simply `forgets' some or all of the structure of an algebraic object is commonly called a forgetful functor (or, an underlying functor). Thus the forgetful functor \(U : \mathbf{Grp} \to \mathbf{Set}\) assigns to each group \(G\) the set \(UG\) of its elements\ldots{}'' (\cite{MacLane1998}, Ch. I §3, p.~14). Brief: \texttt{extends} builds a new structure containing everything an existing one has, plus more, generating a \texttt{.toX} forgetful projection for free.
\end{itemize}

\begin{lstlisting}
structure Point3D extends Point where
  z : Nat

def origin3D : Point3D := { x := 0, y := 0, z := 0 }

#eval origin3D.x   -- inherited field, 0
\end{lstlisting}

\texttt{extends} is used later: a \texttt{CommGroup} (commutative group) will \texttt{extend} \texttt{Group} with one extra axiom (commutativity), instead of repeating all the group fields.

A second single-parent example, extending \texttt{Point} a different way:

\begin{lstlisting}
structure ColorPoint extends Point where
  color : String

def redOrigin : ColorPoint := { x := 0, y := 0, color := "red" }

#eval redOrigin.x       -- inherited field, 0
#eval redOrigin.color   -- own field, "red"
\end{lstlisting}

\subsection{Extending more than one structure at once}\label{sec:02-functions-and-structures:03-extending-structures:3}

\texttt{extends} is not limited to a single parent. \texttt{structure C extends A, B} builds a structure with every field of \texttt{A}, every field of \texttt{B}, and any fields declared in \texttt{C} itself, generating both \texttt{.toA} and \texttt{.toB} forgetful projections:

\begin{lstlisting}
structure Named where
  name : String

structure NamedPoint extends Point, Named where
  color : String

def origin3 : NamedPoint :=
  { x := 0, y := 0, name := "origin", color := "black" }

#eval origin3.x            -- from Point, 0
#eval origin3.name         -- from Named, "origin"
#eval origin3.toPoint.x    -- forgetful projection to Point, 0
#eval origin3.toNamed.name -- forgetful projection to Named, "origin"
\end{lstlisting}

\textbf{Field-name collisions.} If two parents both declare a field with the same name \emph{and} the same type, \texttt{extends} merges them into a single shared field rather than raising an error --- \texttt{NamedPoint} above has no collision, but if both \texttt{Point} and \texttt{Named} had declared an \texttt{x : Nat} field, the resulting structure would have exactly one \texttt{x}, filled once, and readable through either parent's forgetful projection. A collision on the same field name with \emph{different} types, by contrast, is a genuine error (\texttt{Field type mismatch: Field 'x' from parent 'B' has type String but is expected to have type Nat}) --- Lean does not silently pick one type over the other or rename either field.

\begin{mathreading}

\texttt{structure Point3D extends Point where z : Nat} is the product \(\mathrm{Point3D} = \mathrm{Point} \times \mathbb{N} \cong
\mathbb{N} \times \mathbb{N} \times \mathbb{N}\), together with the \emph{forgetful map} \(\mathrm{Point3D} \to \mathrm{Point}\) (projecting away \(z\)), generated automatically as \texttt{.toPoint}. In algebraic language, this is exactly the pattern ``a \(\mathrm{Point3D}\)-structure is a \(\mathrm{Point}\)-structure plus one more piece of data.'' This is precisely how a \texttt{CommGroup} will later be ``a \texttt{Group}-structure plus one more axiom (\(ab = ba\))'': a full subcategory of \texttt{Group} cut out by an extra condition, with the forgetful functor \(\mathrm{CommGroup} \to \mathrm{Group}\) being exactly \texttt{.toGroup}.

\end{mathreading}

\subsection{References}\label{sec:02-functions-and-structures:03-extending-structures:4}

Full citations in the \hyperref[sec:root:bibliography:1]{Bibliography}. Formal definitions are gathered in Recall, above.

\begin{itemize}
\tightlist
\item
  Lean 4 documentation (\cite{LeanDocs}) --- the auto-generated \texttt{.toX} projection produced by \texttt{extends}, used above as \texttt{origin3D.x} and \texttt{.toPoint}.
\item
  Mac Lane (\cite{MacLane1998}), Ch. I §3 ``Functors,'' p.~14 --- forgetful functor.
\end{itemize}

\textbf{Key points.} A \texttt{structure} bundles data (and optionally proofs) under one name, built via the anonymous constructor \texttt{⟨...⟩} and read back out via field projection \texttt{.field}. \texttt{extends} builds a new structure containing everything an existing one has, plus more, generating a \texttt{.toX} forgetful projection for free --- the exact mechanism \texttt{CommGroup} uses to add commutativity to \texttt{Group} in Chapter 6.

\textbf{Socratic questions.}

\begin{enumerate}
\def\labelenumi{\arabic{enumi}.}
\tightlist
\item
  \emph{\texttt{⟨0, 0⟩} and \texttt{\{ x := 0, y := 0 \}} build the identical \texttt{Point}. Since the named form says more, why does the anonymous form ever get used?} Because the \emph{expected type} already says which fields are which --- \texttt{def origin : Point := ⟨0, 0⟩} cannot be ambiguous, since Lean already knows a \texttt{Point} is being built and in what field order. The named form earns its keep only when that context is not enough to make the assignment obvious to a reader.
\item
  \emph{\texttt{Point3D extends Point} generates \texttt{.toPoint} automatically. What would have to be written by hand instead, if \texttt{extends} did not exist?} A separate function \texttt{Point3D.toPoint : Point3D → Point} projecting out the shared fields one at a time --- exactly what a forgetful functor does explicitly, which is the whole reason \texttt{extends} is read categorically rather than as a mere convenience keyword.
\item
  \emph{A \texttt{structure} can bundle proofs as fields, not only data. What changes about }checking* that a term has the right type, once one of its fields is a proof rather than a number?* Nothing about the mechanism --- Lean still checks the field has the stated type --- but the stated type is now a proposition, so supplying that field means supplying a proof, checked once at construction time. This is exactly what makes \texttt{Group} (Chapter 6) impossible to build carelessly: the axiom fields cannot be filled in with nonsense that happens to type-check as data can.
\end{enumerate}

